\subsection{An explicit open elliptic target for the original-factorial examples}
\label{sec:ghost-elliptic-target}

\begin{result}{C.3.1}{Proved}{The pointed Heisenberg cover as an open elliptic curve}
\label{prop:ghost-elliptic-target}
Let \(Y\to U\) be the degree-\(27\) finite etale
cover defined by the class-two exponent-three geometric
quotient \(H_3\) and the point-three section: its
arithmetic fundamental group is the inverse image
of that section in the arithmetic \(H_3\)-quotient.
Its smooth compactification, pointed by the sheet
above three, is
\[
                 E':\ y^2=x^3+9.
\]
The open curve \(Y\) is the complement of the
twenty-seven geometric points above \(0,1,\infty\).
The map to \(U\) is the restriction of the explicit
composite
\[
 E'\xrightarrow{\ \phi\ } E
   \xrightarrow{\ \iota^{-1}\ } C
   \xrightarrow{\ t=3u^3\ }\mathbb P^1,
 \quad
 E:y^2=x^3-243,\quad C:3u^3-2v^3=1,
                                                        \tag{8}
\]
where
\[
 \phi(x,y)=\left(\frac{x^3+36}{x^2},
                    \frac{y(x^3-72)}{x^3}\right)
\]
and, for \(d=3u(u-1)-2v(v-1)\),
\[
 \iota(u,v)=
 \left(\frac{18(uv-1)}d,
       -\frac{27[3u(u+1)+2v(v+1)]}d\right).
                                                        \tag{9}
\]
For every pair constructed in
Theorem~\ref{thm:ghost-original-first-integral-class3},
its ordinary local point at \(q_0\) has no rational
lift to \(Y\). There is no continuous section of the
geometrically pro-three arithmetic fundamental group
of \(Y\) whose pushforward to \(U\) has that
abelian pair.
\end{result}
\begin{proof}
We first identify the curve and the maps, then specify
the section constraint that they impose.

\emph{The cubic quotient and its pointed model.}
Quotienting the pointed \(H_3\)-cover by its center
gives the two Kummer equations
\(t=3u^3\), \(1-t=-2v^3\). Their projective
completion is the smooth diagonal cubic
\(3U^3-2V^3=W^3\), with the distinguished rational
point \(P=(1,1)\) in the affine chart. Formula (9)
satisfies \(y^2=x^3-243\). Here is an explicit
inverse, also verifying that the map has degree one.
For a point \((x,y)\) of \(E\), set
\[
 \begin{gathered}
 r=(y-27)/(3x),\qquad s=(y+27)/(3x),\\
 X=\frac{x^3-972}{9x^2},\qquad
 Y_1=\frac{y(x^3+1944)}{27x^3},\\
 k=\frac{Y_1-15}{X-6},\qquad X_1=k^2-X-6,\\
 u=\frac{X_1s}{6(k-3)},\qquad
 v=\frac{X_1r}{6(k-2)}.
 \end{gathered}                                      \tag{10}
\]
These formulas give \(3u^3-2v^3=1\), and substituting
them in (9) returns \((x,y)\). They are identities
in the function fields, so extend to inverse
isomorphisms of the smooth projective curves. The
second coordinate in (9) has a pole at \(P\),
since its numerator there is nonzero while its
denominator is zero. Thus \(\iota(P)=O_E\).
The retained exact calculation
\path{code/check_elliptic_heisenberg.py}
verifies these identities by polynomial reduction.
It also checks that the tangent third point
\((-1/5,-4/5)\) maps to \((7,-10)\).

For clarity, the elementary degree-three morphism
from this cubic to \(E'\) is
\[
       (u,v)\longmapsto
        (6uv,\,3(3u^3+2v^3));
\]
it sends \(P\) to \((6,15)\), so a translation
is needed before treating it as a pointed isogeny.
Formula (10) includes that translation explicitly.

\emph{Identifying the central cover.}
Each of the three inertia generators of the
geometric \(H_3\)-cover has order three, including
the one above infinity. Riemann--Hurwitz gives
\[
 2g(\overline Y)-2=-2\cdot27
                  +3\cdot27(1-1/3)=0.
\]
The point-three sheet gives a rational origin, so
\(\overline Y\) is an elliptic curve. Its center
quotient \(\overline Y\to C\) has degree three
and is unramified: the center meets each cyclic
inertia group trivially, because the latter has
nonzero image in \(H_3^{\rm ab}\). Via \(\iota\)
this is a pointed isogeny
\(f:\overline Y\to E\).
Origin-preserving maps are group homomorphisms by
\cite[II.1.5--1.6]{GhostMilneEC}.

The displayed map \(\phi:E'\to E\) is a
degree-three isogeny. Its kernel consists of
\(O,(0,3),(0,-3)\). The reverse isogeny
\[
 \psi:E\longrightarrow E',\qquad
 \psi(x,y)=
 \left(\frac{x^3-972}{9x^2},
       \frac{y(x^3+1944)}{27x^3}\right)
\]
satisfies \(\phi\psi=[3]_E\); this identity
is independently checked against the chord-and-tangent
tripling formula in the same program.
The curve \(E\) has exactly one subgroup of order
three defined over \(\mathbb Q\). Indeed, such a
subgroup is \(\{O,Q,-Q\}\), so its common
nonzero \(x\)-coordinate must be rational.
The doubling formula gives its division polynomial
\(3x(x^3-972)\). Since \(972=2^2 3^5\) is
not a rational cube, the only rational root is zero.

Here is the exact comparison with the center quotient.
The kernel of \(f\) is killed by three. The
factorization property for isogenies
\cite[II.1.7--1.8]{GhostMilneEC} therefore factors
\([3]_{\overline Y}\) as
\(\overline Y\xrightarrow f E\xrightarrow g\overline Y\),
where \(g\) has degree three. Surjectivity of
\(f\) gives \(fg=[3]_E\). The kernel of \(g\)
is the unique rational subgroup of order three on
\(E\), hence equals \(\ker\psi\). The same
factorization property gives a degree-one isogeny
\(j:\overline Y\to E'\) with \(jg=\psi\).
Then \(\phi jg=[3]_E=fg\); surjectivity of
\(g\) gives \(\phi j=f\). This proves the
pointed isomorphism and the exact map (8).
There are nine geometric points over each of
\(0,1,\infty\), giving the stated twenty-seven
punctures.

\emph{Which global mod-three section can reach this cover.}
For an arbitrary ordinary pair supplied by A.44
with its original exponents, every integral class-two
section, after its abelian normalization modulo
three, has central class
\[
                       h_3+k\eta\pmod3,
                         \qquad k=0,1,2.             \tag{11}
\]
This is the unramified-cocycle calculation in
Theorem~\ref{thm:ghost-original-first-integral-class3};
it does not use that theorem's extra first-prime
condition. The coefficient-prime-three local
abelian classes at two are exactly those of the
point three. All \(2\)-adic units are uniquely
three-divisible, and the second valuation is one.
Their literal integral cocycles can be aligned
locally by an abelian geometric conjugator divisible
by three: the differences are zero modulo three,
and \(\chi(F_2)-1=1\) is a three-adic unit.
This leaves (11) unchanged modulo three.

At two the central inertia comparison is zero,
and the Frobenius comparison is
\(k\eta(F_2)=k/8\). The second third-residue
formula changes by this value, multiplied by the
second valuation one. The point-three section has
zero third obstruction; therefore the classes
\(k=1,2\) have nonzero local obstruction in
\(\mathbb Z_3/(2^2-1)=\mathbb Z/3\).
Only \(k=0\) can be the reduction of an integral
class-three section.

Consequently every full geometrically pro-three
section above such a pair projects, up to geometric
conjugacy, to the point-three section in the
\(H_3\)-quotient. By the definition of the pointed
cover this is exactly the condition for its image
to lie in the open arithmetic subgroup associated
with \(Y\). Thus it factors through the
geometrically pro-three arithmetic group of \(Y\).
This group is the actual inverse-image subgroup in
the geometrically pro-three arithmetic group of
\(U\): its geometric index is \(27\), and any
open normal subgroup of three-power index in that
geometric subgroup contains its finite intersection
of conjugates, still of three-power index in the
ambient geometric group. Thus the induced pro-three
topology is the full one, and completion identifies
the subgroup rather than merely giving a map to it.
For a class-three section the corresponding assertion
is its factorization through the image of this open
subgroup in the class-three arithmetic quotient;
no lift from that quotient to the full group of
\(Y\) is being inferred. The other two classes
in (11) would define central twists of the cover,
but they have already been excluded at two before
this full-section reduction is made.

\emph{The local fibre exclusion for the prescribed examples.}
At the selected prime \(q_0\) of
Theorem~\ref{thm:ghost-original-first-integral-class3},
the ordinary point \(t_{q_0}\) has valuation
\(24>0\). Its abelian classes modulo three are
both zero: they are the local classes of
\((3,-2)\), and (7) makes both cubic powers.
The local boundary-conjugacy statement
Proposition~\ref{prop:ghost-local-boundary}, applied
to its point section, identifies its full pro-three
image with the cyclic boundary section of its
first Kummer class. That image is trivial in
\(H_3\). The local cyclotomic action modulo
three is also trivial, since \(q_0\equiv1\pmod3\).
By contrast the point-three section has nonzero
central Frobenius in \(H_3\) at \(q_0\), by (7).
These two local sections are not geometrically
conjugate in the finite quotient: conjugation fixes
the center and cannot make its nonzero element zero.

The exact fibre criterion for a pointed regular
geometric cover says that \(Y_{t_{q_0}}\) has a
\(\mathbb Q_{q_0}\)-point precisely when these
two quotient sections are conjugate. Indeed the
geometric fibre is the \(H_3\)-coset set of the
point-three subgroup; a rational point is a fixed
coset under the point section of \(t_{q_0}\).
For the chosen Frobenius the action on this set
is translation by a nonzero central element, so
has no fixed point. This proves the stated local
nonexistence.

Finally any section of the full geometrically
pro-three arithmetic group of \(Y\) pushes forward
continuously to that of \(U\). With the asserted
abelian pair its restriction at \(q_0\) is
geometrically conjugate to the ordinary point
section, by Proposition~\ref{prop:ghost-local-boundary}.
Its finite \(H_3\)-image would therefore have
to be trivial, while factorization through \(Y\)
makes it conjugate to the nontrivial point-three
image. This contradiction proves the section
exclusion using the actual cover map.
\end{proof}

The open elliptic curve is the target in this proof.
The comparison with its compactification has a
precise kernel: the map
\(\pi_1(Y_{\overline{\mathbb Q}})^{(3)}
 \to T_3(E')\) kills the boundary inertia groups.
The pushforward to \(\pi_1(U_{\overline{\mathbb Q}})^{(3)}\)
sends inertia above zero to a conjugate of \(x^3\),
whose first abelian coordinate is the nonzero element
\(3\in\mathbb Z_3(1)\). Hence that pushforward
does not factor through the compact quotient:
its kernel does not contain the compactification
kernel. No assertion here transfers the exclusion
to a section of the compact elliptic curve alone.
